% généré par scripts/06_textes.py
\dosmeta{Année de classement : 384 (formule transmise par l’EDB (lue)). Références : ICUR I 2099 = ILCV 2782 (add). Rattachement selon l’EDB : \emph{Incertae originis} (le volume I des ICVR réunit les inscriptions dont la provenance est perdue). Conservation selon l’EDB : \emph{Musei Vaticani: Galleria Lapidaria}. Support : Plaque de marbre. 64 × 98 cm (h. × l.). Inscription gravée. Fiche EDB : Filippo Piazzolla, première version 2015. Notice du catalogue p.~\pageref{not:268}.}
\txtbloc{Texte transmis (EDCS, état 2022)}{\lnum{1}Hoc tumulum patris filius Benedict[us]\par
\lnum{2}fieri volvit causa amoris paternae re[cordatio]\par
\lnum{3}nis Lucifero homini innocentiss[imo amico(?)]\par
\lnum{4}amicorum qui dum in rebus humanis [fuit]\par
\lnum{5}omnibus paratissimus et dulcis adqu[e]\par
\lnum{6}fuisse videbatur vix enim semper sin[e culpa(?) tem]\par
\lnum{7}poris sui annis LXV et m(ensibus) VI d(iebus) XII decess[it in pace(?)]}
\txtbloc{Texte de l’EDB (EDB32083)}{\lnum{1}hoc tumulum patris filius̲ B̲e̲n̲e̲d̲i̲c̲t̲[us]\par
\lnum{2}fieri voluit causa amoris pa̲t̲e̲r̲n̲a̲e̲ r̲e̲[cordatio]/nis Lucifero homini innoce̲n̲t̲i̲s̲s̲[imo amico?]\par
\lnum{3}amicorum qui dum in rebus hu̲m̲a̲n̲i̲s̲ [fuit]\par
\lnum{4}omnibus paratissimus et dulc̲i̲s̲ a̲d̲q̲u̲̲[e]\par
\lnum{5}fuisse videbatur vix enim semp̲e̲r̲ s̲i̲n̲[e culpa? tem]/poris sui annis LXV et m(ensibus) VI d(iebus) XI̲I̲ d̲e̲c̲e̲s̲s̲[it in pace?]\par
\lnum{6}VII k(alendas) iul(ias) Ricomere et Clearcho cons(ulibu)s ((palma)) ((avis))}
